% Part: first-order-logic
% Chapter: tableaux
% Section: proof-theoretic-notions

\documentclass[../../../include/open-logic-section]{subfiles}

\begin{document}

\iftag{FOL}
      {\olfileid[hi]{fol}{tab}{ptn}}
      {\olfileid[hi]{pl}{tab}{ptn}}
      
\olsection{प्रमाण-सैद्धांतिक धारणाएँ}

\begin{editorial}
यह अनुभाग सत्य-वृक्षों के लिए सिद्धता-संबंध और संगतता की परिभाषाएँ एकत्र
करता है।
\end{editorial}

\begin{explain}
जिस प्रकार हमने कई महत्त्वपूर्ण अर्थवैज्ञानिक धारणाएँ (वैधता, तार्किक
निष्पत्ति, संतुष्टनीयता) परिभाषित की हैं, उसी प्रकार अब उनकी संगत
\emph{प्रमाण-सैद्धांतिक धारणाएँ} परिभाषित करते हैं। इनकी परिभाषा में यह
सहारा नहीं लिया जाता कि !!{sentence}s को कौन-सी !!{structure}s संतुष्ट करती
हैं, बल्कि कुछ बंद !!{tableau}ों के अस्तित्व का सहारा लिया जाता है। इन
धारणाओं का एक-दूसरे से मेल खाना एक महत्त्वपूर्ण खोज थी। यह मेल
\emph{सुदृढ़ता} और \emph{पूर्णता प्रमेयों} का विषय है।
\end{explain}


\begin{defn}[प्रमेय]
!!{sentence}~$!A$ \emph{प्रमेय} है यदि
$\sFmla{\False}{!A}$ के लिए कोई बंद !!{tableau} हो। यदि ऐसा हो तो
$\Proves !A$ लिखते हैं; यदि $!A$ प्रमेय नहीं है तो $\Proves/ !A$ लिखते हैं।
\end{defn}

\begin{defn}[!!^{derivability}]
!!^a{sentence}~$!A$ किसी समुच्चय से \emph{!!{derivable}} है; यदि वह
!!{sentence}s का समुच्चय~$\Gamma$ हो, तो इसे $\Gamma \Proves !A$ लिखते हैं।
यह तभी और केवल तभी है जब कोई परिमित समुच्चय
$\{!B_1, \dots, !B_n\} \subseteq \Gamma$ हो और निम्न समुच्चय के लिए कोई
बंद !!{tableau} हो:
\[
\{
  \sFmla{\False}{!A},
  \sFmla{\True}{!B_1}, \dots,
  \sFmla{\True}{!B_n}
  \}.
\]
यदि $!A$, $\Gamma$ से !!{derivable} नहीं है तो $\Gamma \Proves/
!A$ लिखते हैं।
\end{defn}

\begin{defn}[संगतता]
!!{sentence}s का समुच्चय~$\Gamma$ \emph{असंगत} है तभी और केवल तभी जब कोई
परिमित समुच्चय $\{!B_1, \dots, !B_n\} \subseteq \Gamma$ हो और निम्न
समुच्चय के लिए कोई बंद !!{tableau} हो:
\[
\{
  \sFmla{\True}{!B_1}, \dots,
  \sFmla{\True}{!B_n}  
  \}.
\]
यदि $\Gamma$ असंगत नहीं है, तो उसे \emph{संगत} कहते हैं।
\end{defn}

\begin{prop}[स्वतुल्यता]
\ollabel{prop:reflexivity}
यदि $!A \in \Gamma$, तो $\Gamma \Proves !A$।
\end{prop}

\begin{proof}
यदि $!A \in \Gamma$, तो $\{!A\}$,~$\Gamma$ का परिमित उपसमुच्चय है और
निम्न !!{tableau} बंद है:
\begin{oltableau}
  [\sFmla{\False}{\formula{A}}, just = \TAss
    [\sFmla{\True}{\formula{A}}, just = \TAss,close]
  ]
\end{oltableau}
\end{proof}
  

\begin{prop}[एकदिष्टता]
\ollabel{prop:monotonicity}
यदि $\Gamma \subseteq \Delta$ और $\Gamma \Proves !A$, तो $\Delta
\Proves !A$।
\end{prop}

\begin{proof}
$\Gamma$ का प्रत्येक परिमित उपसमुच्चय,~$\Delta$ का भी परिमित उपसमुच्चय है।
\end{proof}

\begin{prop}[संक्रामकता]
\ollabel{prop:transitivity}
यदि $\Gamma \Proves !A$ और $\{!A\} \cup
\Delta \Proves !B$, तो $\Gamma \cup \Delta \Proves !B$।
\end{prop}

\begin{proof}
यदि $\{!A\} \cup \Delta \Proves !B$, तो कोई परिमित उपसमुच्चय $\Delta_0 =
\{!C_1, \dots, !C_n\} \subseteq \Delta$ ऐसा है कि
\begin{align*}
\{\sFmla{\False}{!B}, & \sFmla{\True}{!A}, \sFmla{\True}{!C_1},
\dots, \sFmla{\True}{!C_n}\}
\intertext{के लिए कोई बंद !!{tableau} है। यदि $\Gamma \Proves !A$, तो
  $!D_1$, \dots, $!D_m \subseteq \Gamma$ ऐसे हैं कि}
\{\sFmla{\False}{!A}, & \sFmla{\True}{!D_1},
\dots, \sFmla{\True}{!D_m}\}
\end{align*}
के लिए कोई बंद !!{tableau} है।

अब निम्न मान्यताओं वाले !!{tableau} पर विचार करें:
\[
\sFmla{\False}{!B},
\sFmla{\True}{!C_1}, \dots, \sFmla{\True}{!C_n},
\sFmla{\True}{!D_1}, \dots, \sFmla{\True}{!D_m}.
\]
$!A$ पर \Cut{} नियम लगाएँ। इससे दो शाखाएँ बनती हैं: एक में
$\sFmla{\True}{!A}$ और दूसरी में $\sFmla{\False}{!A}$ है। अतः पहली शाखा
में निम्न सभी उपलब्ध हैं:
\[
\{\sFmla{\False}{!B}, \sFmla{\True}{!A},
\sFmla{\True}{!C_1}, \dots, \sFmla{\True}{!C_n}\}
\]
इन मान्यताओं के लिए बंद !!{tableau} होने के कारण हम उसे इस शाखा से जोड़
सकते हैं; $\sFmla{\True}{!A}$ से होकर जाने वाली प्रत्येक शाखा बंद हो जाती
है। दूसरी शाखा में निम्न सभी उपलब्ध हैं:
\[
\{\sFmla{\False}{!A}, \sFmla{\True}{!D_1}, \dots,
\sFmla{\True}{!D_m}\}
\]
अतः दूसरी ओर को भी पूरा करके बंद !!{tableau} पाया जा सकता है। इससे
$\Gamma \cup \Delta \Proves !B$ सिद्ध होता है।
\end{proof}

ध्यान दें कि विशेषतः इसका अर्थ है: यदि $\Gamma \Proves !A$ और $!A
\Proves !B$, तो $\Gamma \Proves !B$। इससे यह भी मिलता है कि यदि $!A_1,
\dots, !A_n \Proves !B$ और $\Gamma \Proves !A_i$ प्रत्येक~$i$ के लिए, तो
$\Gamma \Proves !B$।

\begin{prop}
\ollabel{prop:incons}
$\Gamma$ असंगत है तभी और केवल तभी जब $\Gamma \Proves !A$ प्रत्येक
!!{sentence}~$!A$ के लिए हो।
\end{prop}

\begin{proof}
अभ्यास।
\end{proof}

\tagprob{FOL}
\begin{prob}
\olref[fol][tab][ptn]{prop:incons} सिद्ध कीजिए।
\end{prob}
\tagendprob

\tagprob{notFOL}
\begin{prob}
\olref[pl][tab][ptn]{prop:incons} सिद्ध कीजिए।
\end{prob}
\tagendprob

\begin{prop}[संहतता]
  \ollabel{prop:proves-compact}
  \begin{enumerate}
  \item यदि $\Gamma \Proves !A$, तो कोई परिमित उपसमुच्चय $\Gamma_0
    \subseteq \Gamma$ ऐसा है कि $\Gamma_0 \Proves !A$।
  \item यदि~$\Gamma$ का हर परिमित उपसमुच्चय संगत है, तो $\Gamma$ संगत है।
  \end{enumerate}
\end{prop}

\begin{proof}
  \begin{enumerate}
    \item यदि $\Gamma \Proves !A$, तो कोई परिमित उपसमुच्चय
      $\Gamma_0 = \{!B_1, \dots, !B_n\}$ और निम्न के लिए बंद !!{tableau}
      है:
      \[
      \{\sFmla{\False}{!A}, \sFmla{\True}{!B_1}, \dots, \sFmla{\True}{!B_n}\}
      \]
      यही !!{tableau}, $\Gamma_0 \Proves !A$ भी दिखाता है।
    \item यदि $\Gamma$ असंगत है, तो किसी परिमित उपसमुच्चय
      $\Gamma_0 = \{!B_1, \dots, !B_n\}$ के लिए निम्न बंद !!{tableau} है:
      \[
      \{\sFmla{\True}{!B_1}, \dots, \sFmla{\True}{!B_n}\}
      \]
      यह बंद !!{tableau} दिखाता है कि $\Gamma_0$ असंगत है।
  \end{enumerate}
\end{proof}

\end{document}
